%(BEGIN_QUESTION)
% Copyright 2006, Tony R. Kuphaldt, released under the Creative Commons Attribution License (v 1.0)
% This means you may do almost anything with this work of mine, so long as you give me proper credit

En jernblokk og en kobberblokk varmes begge opp til 200$^{o}$ F i en ovn, og legges deretter på et trebord for å kjøles ned til romtemperatur. Hvis alle andre faktorer er like (masse, overflateareal osv.), hvilken blokk vil kjøles ned raskest? Hvorfor?

\underbar{file i00330}
%(END_QUESTION)





%(BEGIN_ANSWER)

Kobberblokken vil kjøles ned raskere, fordi kobber har lavere spesifikk varmekapasitet enn jern. Det betyr at en gitt mengde avgitt varmeenergi senker temperaturen mer i kobber enn i jern ved samme energitap. En annen måte å si det på er at kobberblokken inneholder mindre varmeenergi enn jernblokken, selv om de starter ved samme temperatur.

\vskip 10pt

$c_{Cu}$ = 0.093 cal/g$\cdot ^{o}$C

\vskip 10pt

$c_{Fe}$ = 0.113 cal/g$\cdot ^{o}$C

%(END_ANSWER)





%(BEGIN_NOTES)

Materialet i bordet (tre) er overflødig informasjon, inkludert for å utfordre studenten til å vurdere om opplysningene er relevante for å løse oppgaven.

%INDEX% Physics, heat and temperature: calorimetry problem 

%(END_NOTES)
